\section{Results}
\label{sec:results}

\subsection{Full-array subject-independent performance}
With all 16 electrodes, XGBoost reached a Macro-F1 of \FullXGB\ (95\% CI \FullXGBLCB--\FullXGBUCB) over the 43 held-out participants. The support vector machine reached \FullSVM\ (\FullSVMLCB--\FullSVMUCB) and the random forest \FullRF\ (\FullRFLCB--\FullRFUCB) (table~\ref{tab:full_channel}, figure~\ref{fig:full}). No statistically detectable difference was observed between XGBoost and the support vector machine (paired difference $\XGBvsSVMDiff$, 95\% CI $\XGBvsSVMLCB$ to $\XGBvsSVMUCB$, $p_{\mathrm{Holm}} \XGBvsSVMP$, $r_{rb} = \XGBvsSVMR$). Both exceeded the random forest ($p_{\mathrm{Holm}} \XGBvsRFP$ and $\SVMvsRFP$; $r_{rb} = \XGBvsRFR$ and $\SVMvsRFR$). The spread between participants was far larger than any difference between models: XGBoost ranged from \FullXGBMin\ to \FullXGBMax\ across participants.

Participant-level performance for every active gesture is reported in the supplementary material. These summaries were reconstructed directly from the saved full-array predictions.

\input{tables/tab_full_channel.tex}

\begin{figure}[t]
\centering
\includegraphics[width=\textwidth]{figures/fig3_full_channel.pdf}
\caption{Full-array leave-one-subject-out performance on GRABMyo (43 participants). (a) Macro-F1 of three classifiers; each point is one held-out participant, grey lines join the same participant, black bars mark medians and diamonds mark means with 95\% cluster-bootstrap confidence intervals. The support vector machine and random forest were trained on samples of 40,000 and 200,000 windows. (b) Paired XGBoost minus support vector machine difference per participant, sorted; the grey band is the 95\% confidence interval of the mean difference.}
\label{fig:full}
\end{figure}

\subsection{H1: a 12-electrode subset is non-inferior within the prespecified margin}
Performance declined smoothly as electrodes were removed (figure~\ref{fig:curve}, table~\ref{tab:sparse_ni}). The reference of the curve ($k = 16$) scored \CurveKSixteen, slightly below the full-data result because all curve models share one 200,000-window training sample. At $\delta = 0.02$, the smallest non-inferior subset was $k^* = \KStar$ electrodes ($\bar\Delta_{12} = \DeltaKTwelve$, 95\% CI $\DeltaKTwelveLCB$ to $\DeltaKTwelveUCB$). At $k = 11$ the lower bound ($\DeltaKElevenLCB$) crossed the margin. Every $k$ from 12 to 15 was non-inferior, so the result is not an isolated crossing. With Holm correction over the 15 subset sizes, $k^*$ was unchanged at all three margins ($k = 12$: $p_{\mathrm{Holm}} \HolmPKTwelve$; supplementary table~S1). The sensitivity margins gave $k^* = \KStarMarginPtOne$ at $\delta = 0.01$ and $k^* = \KStarMarginPtThree$ at $\delta = 0.03$. The 12-electrode subset was therefore non-inferior within the prespecified engineering tolerance, a reduction of \KStarReductionPct\% of the array. Larger reductions were not: halving the array to $k = 8$ cost $\DeltaKEight$ (95\% CI $\DeltaKEightLCB$ to $\DeltaKEightUCB$), and $k = 4$ cost $\DeltaKFour$. All paired differences at $k < 16$ were negative. The electrodes most often excluded at $k = 12$ were \DroppedFourSet\ (\DroppedFourFolds\ of 43 folds).

\input{tables/tab_sparse_ni.tex}

\begin{figure}[t]
\centering
\includegraphics[width=\textwidth]{figures/fig4_sparse_curve.pdf}
\caption{Electrode-count performance curve on GRABMyo (43 participants). (a) Mean Macro-F1 of models retrained on the top-$k$ electrodes of each fold's cross-fitted ranking, for the four values of the stability weight $\lambda$ (the curves overlap); the band is the 95\% cluster-bootstrap confidence interval for $\lambda = 0$. (b) Paired difference to the $k = 16$ model with 95\% confidence interval. The thick dashed line marks the primary margin ($-\delta$, $\delta = 0.02$); the thin dotted and dash-dotted grey lines mark the sensitivity margins (0.01 and 0.03).}
\label{fig:curve}
\end{figure}

\subsection{H2: attribution-ranked subsets exceed the mean of random subsets}
Under the shared comparison model (table~\ref{tab:baselines}, figure~\ref{fig:baselines}), SR-CFA subsets outperformed the mean of random subsets at every tested $k$ (Holm-corrected, 8 tests). The advantage was largest at intermediate sizes: $\HTwoKSixDiff$ Macro-F1 at $k = 6$ (95\% CI $\HTwoKSixLCB$ to $\HTwoKSixUCB$, $p_{\mathrm{Holm}} \HTwoKSixP$, $r_{rb} = \HTwoKSixR$) and $\HTwoKEightDiff$ at $k = 8$. It was small at $k = 4$ ($\HTwoKFourDiff$) and at $k = 12$ ($\HTwoKTwelveDiff$, $p_{\mathrm{Holm}} \HTwoKTwelveP$). Within each participant's own random subsets, the SR-CFA subset had a median percentile of \PctKFour\ at $k = 4$, \PctKSix\ at $k = 6$, \PctKEight\ at $k = 8$ and \PctKTwelve\ at $k = 12$, and it reached the 95th percentile for \PctKAboveNinetyFiveSix\ of 43 participants at $k = 6$ and \PctKAboveNinetyFiveTwelve\ at $k = 12$. Attribution-ranked subsets therefore exceeded the average random subset at every size tested but seldom reached the upper tail of a participant's random subsets.

In exploratory, uncorrected contrasts, no statistically detectable difference between SR-CFA and permutation importance was observed at $k = 4$, 8 or 12 (for example, $k = 8$: $\VSPermKEightDiff$, $p \VSPermKEightP$); at $k = 6$, SR-CFA was higher by $\VSPermKSixDiff$ ($p \VSPermKSixP$). Likewise, no statistically detectable difference between SR-CFA and geometric selection, which spreads electrodes evenly around both rings without using data, was observed at $k = 4$ ($\VSGeomKFourDiff$, $p \VSGeomKFourP$), $k = 8$ ($\VSGeomKEightDiff$, $p \VSGeomKEightP$) or $k = 12$ ($\VSGeomKTwelveDiff$, $p \VSGeomKTwelveP$); at $k = 6$, SR-CFA was higher by $\VSGeomKSixDiff$ ($p \VSGeomKSixP$). Mutual information gave the lowest scores at every $k \le 8$.

\input{tables/tab_baselines.tex}

\begin{figure}[t]
\centering
\includegraphics[width=\textwidth]{figures/fig5_selection_baselines.pdf}
\caption{Selection methods at matched subset size $k$ (GRABMyo, 43 participants; one comparison model shared by all methods). (a) Mean Macro-F1 with 95\% cluster-bootstrap confidence intervals, grouped by method type: attribution rankings (stability-regularised cross-fitted attribution, SR-CFA, and cross-fitted SHAP), model-based rankings (permutation importance and mutual information), data-free geometric selection and the mean over random subsets. (b) Percentile of the SR-CFA subset within each participant's own random subsets (one point per participant; bar = median; 100 random subsets per participant, 500 at $k = 12$).}
\label{fig:baselines}
\end{figure}

\subsection{H3: the stability penalty produced no measurable improvement}
The instability term was small relative to attribution magnitude: the median $U_c/A_c$ was \UOverAMedian\% (maximum \UOverAMax\%). Penalising it therefore rarely changed a ranking. Relative to $\lambda = 0$, the top-12 set changed in \TopTwelveDiffersLambdaPtTwoFive, \TopTwelveDiffersLambdaPtFive\ and \TopTwelveDiffersLambdaOne\ of 43 folds at $\lambda = 0.25$, 0.5 and 1, and across all four $k$ tested at least \HThreeMinTied\ of 43 participants received the identical subset under both rankings. None of the \HThreeNTests\ contrasts was significant after Holm correction (largest mean difference \HThreeMaxAbsDiff; smallest $p_{\mathrm{Holm}} \HThreeMinP$; supplementary table~S2). The four curves in figure~\ref{fig:curve}(a) differ by at most \LambdaCurveMaxSpread, and $k^*$ was identical under every $\lambda$ and margin. H3 was not supported.

\subsection{H4: deletion with retraining supports faithfulness}
The expected ordering, Remove-Top $<$ Remove-Random $<$ Remove-Bottom, held at every number of electrodes retained (table~\ref{tab:faithfulness}, figure~\ref{fig:faithfulness}); \FaithNSig\ of the \FaithNTests\ Holm-corrected contrasts were significant. With six electrodes retained, Remove-Top scored \FaithTopKSix, Remove-Random \FaithRandomKSix\ and Remove-Bottom \FaithBottomKSix\ (Remove-Top minus Remove-Random $\FaithTRKSixDiff$, $p_{\mathrm{Holm}} \FaithTRKSixP$; Remove-Random minus Remove-Bottom $\FaithRBKSixDiff$, $p_{\mathrm{Holm}} \FaithRBKSixP$). The only non-significant contrast was Remove-Top versus Remove-Bottom with 14 electrodes retained ($\FaithTBKFourteenDiff$, $p_{\mathrm{Holm}} \FaithTBKFourteenP$), where only two electrodes are removed. The separation between deletion orders was largest at intermediate sizes and small at the extremes.

\input{tables/tab_faithfulness.tex}

\begin{figure}[t]
\centering
\includegraphics[width=0.62\textwidth]{figures/fig6_faithfulness.pdf}
\caption{Deletion-with-retraining test of attribution faithfulness on GRABMyo (43 participants; mean and 95\% cluster-bootstrap confidence interval). Remove-Top deletes the highest-ranked electrodes, Remove-Bottom the lowest-ranked and Remove-Random a random set; all use the same comparison model and training sample. Remove-Random is shown where random subsets were sampled (4, 6, 8 and 12 electrodes retained).}
\label{fig:faithfulness}
\end{figure}

\subsection{Stability of the ranking across participants}
Attribution was concentrated and reproducible (figure~\ref{fig:attribution}). The same four electrodes (\TopFourSet) formed the top four in \TopFourFolds\ of 43 outer folds and carried \TopFourShare\% of the summed attribution, whereas the eight lowest-ranked electrodes carried \BottomEightShare\%. Rankings computed from individual participants' attributions agreed closely (table~\ref{tab:stability}): mean Spearman $\rho$ = \StabSubjRho\ $\pm$ \StabSubjRhoSD\ over \StabSubjPairs\ pairs, and a mean Jaccard index of the top six of \StabSubjJacSix. Rankings from disjoint inner folds agreed at $\rho$ = \StabInnerRho. The outer-fold agreement ($\rho$ = \StabOuterRho) is descriptive rather than an independent stability estimate, because those training sets share about 98\% of their participants. Across folds, the highest- and lowest-ranked electrodes kept their positions, and the rank ranges were widest in the middle of the ordering (figure~\ref{fig:attribution}(b)).

\input{tables/tab_stability.tex}

\begin{figure}[t]
\centering
\includegraphics[width=\textwidth]{figures/fig7_attribution.pdf}
\caption{Cross-fitted electrode attribution on GRABMyo ($\lambda = 0$, 43 outer folds). (a) Mean share of attribution per electrode, drawn on the $2 \times 8$ montage; each ring~2 electrode is aligned with, and 2~cm from, the ring~1 electrode at the same circumferential position. (b) Rank of each electrode across inner folds (median and range).}
\label{fig:attribution}
\end{figure}

\subsection{Replication on independent datasets}
\paragraph{NinaPro DB5 (exploratory)} With all 16 channels, Macro-F1 was \DBFiveFull\ $\pm$ \DBFiveFullSD\ (12 movements; chance about 0.083), far below GRABMyo and with wide between-subject spread (\DBFiveFullMin--\DBFiveFullMax) (figure~\ref{fig:generalisation}(b), supplementary material). At $\delta = 0.02$, $k = 12$ met the margin by a negligible amount ($\bar\Delta_{12} = \DBFiveDeltaKTwelve$, lower bound $\DBFiveDeltaKTwelveLCBSix$): it was non-inferior under \DBFiveSeedFrac\% of \DBFiveSeedN\ alternative bootstrap seeds, $k = 13$, 14 and 15 all failed (lower bounds $\DBFiveDeltaKThirteenLCB$, $\DBFiveDeltaKFourteenLCB$ and $\DBFiveDeltaKFifteenLCB$), and with Holm correction no subset was non-inferior. With ten subjects the estimates have wide uncertainty, and DB5 neither supports nor contradicts the GRABMyo reduction; its mean curve increased broadly with $k$ but was not monotone.

\begin{figure}[t]
\centering
\includegraphics[width=\textwidth]{figures/fig8_generalisation.pdf}
\caption{Exploratory analyses. (a) GRABMyo cross-session evaluation: each point is one held-out participant, grey lines join the same participant across conditions and bars mark medians; every condition is also subject-independent. (b--c) Mean paired difference to the full array of the same curve, $\bar\Delta_k$, with 95\% cluster-bootstrap confidence interval, for (b) NinaPro DB5 (10 subjects, reference $k = 16$), (c) GRABMyo derived bipolar pairs (43 participants, reference $k = 8$). The dashed line marks the primary margin; filled markers indicate a lower bound above the margin and open markers a lower bound at or below it.}
\label{fig:generalisation}
\end{figure}

\input{generated/participant_extension_results.tex}
\begin{figure}[t]
\centering
\includegraphics[width=\textwidth]{figures/fig9_participant_extension.pdf}
\caption{Post-hoc participant-level attribution extension. Every cell is the mean paired difference across the GRABMyo participants, comparing person-level with fold-level aggregation under the same penalty, subset size and model protocol. Blue favours person-level aggregation and orange favours fold aggregation. Asterisks identify Holm-corrected differences within the complete family. All candidates are shown; outer-test outcomes did not select a penalty.}
\label{fig:participant_extension}
\end{figure}

\subsection{Other exploratory analyses}
\paragraph{Cross-session} Training on session~1 and testing on sessions 1, 2 and 3 gave \CSOneSOne, \CSOneSTwo\ and \CSOneSThree\ (supplementary material, figure~\ref{fig:generalisation}(a)). Mean Macro-F1 on sessions 2 and 3 was numerically higher than on session 1; the session-2 difference was nominally significant ($\CSSTwoVsSOneDiff$, 95\% CI $\CSSTwoVsSOneLCB$ to $\CSSTwoVsSOneUCB$, uncorrected $p \CSSTwoVsSOneP$), and no statistically detectable difference was observed for session 3 ($\CSSThreeVsSOneDiff$, $p \CSSThreeVsSOneP$). Because every condition is also subject-independent, this does not show that sessions are interchangeable, only that no additional cross-day loss was detectable on top of the cross-user gap. Training on sessions 1 and 2 improved session-3 performance over training on session 1 alone ($\CSTwoVsOneDiff$, 95\% CI $\CSTwoVsOneLCB$ to $\CSTwoVsOneUCB$, $p \CSTwoVsOneP$), but this design cannot separate day diversity from the doubled amount of training data.

\input{generated/bipolar_results.tex}
